\section{Methodology}
To address the research question, CLAIM evaluations generated by three LLMs were compared with those provided by human reviewers. Those were obtained from a previously published systematic review, which included 209 studies of AI cardiac MRI segmentation. Each study was evaluated using the 42 items of the 2020 version of CLAIM and organised into four domains: Study description, Dataset description, Model description and Model performance (Appendix \ref{app:table}) \cite{alabed2022quality}. It provided descriptive characteristics for each paper, together with compliance scores for each domain. Therefore, all comparisons in this study were performed using domain-level and overall scores rather than individual CLAIM items. The human scores were used as the reference standard, although they may also involve some degree of subjectivity, as acknowledged by the original authors \cite{alabed2022quality}. 

For this meta-research, the first 60 papers in the original dataset with available full texts were selected. Full-text access was obtained either through publicly accessible links or through the institutional access provided by Universitat Pompeu Fabra. 

\subsection{LLMs assessments and compliance score calculation}
Three of the most widely used LLMs were included in the study: Gemini 3.8 Flash (Google), GPT-6 (OpenAI) and Claude Sonnet 5.5 (Anthropic). A new independent Google account was created to avoid any influence from previous conversations, saved information or custom instructions. All models were used with their standard reasoning mode, without enabling extended-reasoning options or changing the default settings. 

The evaluations were generated using the same prompt (Appendix \ref{app:prompt}) for all three models and all 60 papers. It instructed each model to act as an experienced peer reviewer with expertise in artificial intelligence and medical imaging and to evaluate the manuscript according to all 42 items of CLAIM 2020. The full description of each item was included in the prompt. 

For every item, the models were required to provide one of three possible answers: Yes, No or Not Applicable. They were instructed to rely exclusively on information explicitly reported in the manuscript, assess each item independently and avoid assuming that unreported procedures had been performed. When relevant information was missing, the answer had to be No rather than Not Applicable. To ensure consistent and easily comparable outputs, the responses were restricted to 43 lines: the paper’s PDF filename followed by one answer for each of the 42 CLAIM items, without additional explanations.

Each PDF was uploaded together with the complete prompt between November 6th and 9th of 2026. Manuscripts were processed within the same conversation until the usage limit was reached. Processing was then paused until the limit was reset. 

The responses obtained from each model were transferred to an Excel spreadsheet, with one row per paper evaluation and one column per CLAIM item. This resulted in an expanded dataset of 60 papers for each LLM (240 observations). Based on these responses, a compliance score (Equation \ref{eq:compliance}) was calculated for each paper, model and CLAIM domain, as well as for the overall checklist. 

\begin{equation}
\text{Compliance Score} = \frac{\text{Yes}}{\text{Yes}+\text{No}} \times 100
\label{eq:compliance}
\end{equation}

\subsection{Evaluation metrics}
All statistical analyses and visualizations were performed in Python (version 3.13.18) using Jupyter Notebooks. The data were processed using the pandas and NumPy libraries, statistical analyses were conducted using SciPy and statsmodels, and figures were generated using Matplotlib and Seaborn.

%All statistical analyses and visualizations were performed in Python (version 3.13.18) using Jupyter Notebooks. The data were processed using the pandas (version 2.2.3) and NumPy (version 2.4.6) libraries, statistical analyses were conducted using SciPy (version 1.16.3) and statsmodels (version 0.15.0), and figures were generated using Matplotlib (version 3.10.0) and Seaborn (version 0.13.2).

Statistical analyses were performed separately for the four CLAIM domain scores and the overall compliance score, resulting in five outcomes. For each paper, model and outcome, the error was defined as the difference between the LLM score and the corresponding human score, expressed in percentage points.

Before conducting the statistical analyses, scores were prepared for comparison. Tied values were assigned average ranks in the Friedman and Spearman analyses, while zero differences were excluded from the Wilcoxon signed-rank tests. An adjusted p-value below 0.05 was considered statistically significant for all tests. 

To measure agreement with the human reviewers, the mean absolute error (MAE) was calculated for each model and outcome. To assess whether the models tended to be stricter or more lenient than the human reviewers, the mean error (ME), calculated as the average of the signed differences, was used. Negative ME values indicate a stricter evaluation and positive values a more lenient one. Whether these differences were statistically significant was tested using a two-sided Wilcoxon signed-rank test on the paired scores. The resulting 15 p-values were adjusted using the Holm method. 

To compare the performance of the three LLMs, the absolute errors relative to the human scores were compared across papers using the Friedman test, separately for each of the five outcomes. Again, the five resulting p-values were adjusted using the Holm method. In addition, pairwise comparisons between the three models were performed for each outcome using two-sided Wilcoxon signed-rank tests on the paired absolute errors, regardless of the Friedman test results. P-values from these pairwise comparisons were adjusted using the Holm method separately for each outcome.

Finally, Spearman’s rank correlation coefficient $(\rho)$ was calculated to examine the relationship between the scores. Correlations were assessed between each LLM and the human reviewers, as well as between each pair of LLMs, for all five outcomes. Spearman’s correlation indicates whether papers are ranked similarly. The p-values from the 15 correlations between LLMs and human reviewers and the 15 correlations between pairs of LLMs were adjusted separately using the Holm method. 

%si queremos poner lo del metadatos descomentar los siguiente: 

%A dditionally, the relationship between evaluation scores and paper metadata was explored to assess whether LLMs showed similar evaluation patterns to human assessors across different types of papers. Numerical and categorical variables were visualised using scatter plots and strip plots, respectively.